\paragraph{Proposition E.2.18 (Proved: the cusp quadratic attachment).}
Keep the right-module presentation of E.2.16 and write
\[
 R=\mathbb ZK,\qquad M=\pi_2(E)=eR\oplus T,
 \qquad T=\overline R^{\,2}\oplus L.
\]
For an abelian group $A$, let $\Gamma(A)$ be Whitehead's universal
quadratic group, with quadratic generator $q(a)$ and symmetric
bilinear cross effect
$[a,b]=q(a+b)-q(a)-q(b)$.  In particular $[a,a]=2q(a)$.
There is a lift of the normalized cusp component
$S=\mathbb {CP}^2\#3\overline{\mathbb {CP}}^{\,2}$ to the
universal cover $\widetilde E$ whose fundamental class has image
\[
 \pm\bigl([e(h-1),t]-12q(t)\bigr)
       \quad\hbox{under}\quad
 H_4(\widetilde E;\mathbb Z)\longrightarrow\Gamma(M).
 \tag{E.Q1}
\]
The sign is solely the orientation chosen on the normalized
component.  Its mixed component in
\[
 \Gamma(M)=\Gamma(eR)\oplus(eR\otimes_{\mathbb Z}T)
                                      \oplus\Gamma(T)
 \tag{E.Q2}
\]
is nonzero.  Consequently $E$ is not homotopy equivalent to
$E_0\#(S^2\times S^4)$ for any connected compact six-manifold
$E_0$, even without requiring the equivalence to preserve the
boundary.  Thus the free $R$ summand of its second homotopy group
does not come from a sphere-product connected summand.
Likewise the actual nine-dimensional fibre $Y$ is not
homotopy equivalent to $Y_0\#(S^2\times S^7)$ for any
connected closed nine-manifold $Y_0$, despite the twenty-four
free $\mathbb Z\Pi$ summands in (E.P6).

\begin{proof}
We use the natural, deck-equivariant part of Whitehead's exact
sequence
\[
 H_4(\widetilde E)\xrightarrow{\omega_E}\Gamma(\pi_2(E))
       \longrightarrow\pi_3(E)\longrightarrow
 H_3(\widetilde E)\longrightarrow0.
 \tag{E.Q3}
\]
The exactness and identification of its quadratic term are
\cite[Theorems 1 and 5]{mapsWhitehead50}; the direct-sum rule
(E.Q2) is proved there in Section~2.  These statements apply
to the simply connected universal cover, with no finiteness
requirement.  Naturality, rather than a choice of a word for
one attaching cell, will specify the element in (E.Q1).

For the normalization calculation take the negatives of
the complex-oriented Picard homology generators and denote
them by $(H,E_1,E_2,E_3)$.  Their intersection matrix is
still $\operatorname{diag}(1,-1,-1,-1)$.  This common reversal
matches the phase-cell orientations of (B.5.16): the toric
line with class $-(uw+6\gamma\delta)$ has complex curve
degrees $(-6,0,1)$, whereas its values on the corresponding
phase classes are $(6,0,-1)$; the support function $rs$
similarly fixes the middle sign.  The exact wall-degree
calculation in D.1.9 gives this comparison.  With these
orientations the six conductor classes are
\[
 E_1,\ H-E_1-E_2,\ E_2,\ H-E_2-E_3,
                                      E_3,\ H-E_1-E_3.
\]
The first opposite pair is self-glued in the surviving $u$
direction.  In the gamma cover its relation is therefore
$E_1=H-E_2-E_3$.  The other two pairs connect adjacent gamma
translates.  Their relations, after eliminating $H$, are
$(h-1)E_2=(h-1)E_3=0$.  Thus the image of normalization
homology is, with these actual curve generators,
\[
 A E_1\oplus(A/(h-1))E_2\oplus(A/(h-1))E_3,
 \qquad A=\mathbb Z[h,h^{-1}],\qquad H=E_1+E_2+E_3.
 \tag{E.Q4}
\]
This also follows directly from the transpose of the matrix
$B(h)$ of E.2.12, with the intersection signs in the Picard
basis retained.  It is the whole $H_2$ of the gamma cover:
the degree-zero normalization row is the cellular complex
of an infinite cylinder, whose second homology is zero.

We fix the comparison with the cusp basis of (E.C12), rather
than identify its free summand only abstractly.  In the ordinary
central fibre the first conductor class $\overline C_1$ is the
self-$u$ pair; the other two are the images of $E_2,E_3$,
in either order.  The chosen cusp basis is
\[
 E=-\overline C_1-\overline C_2,\qquad
 U=\overline C_2+\overline C_3,\qquad D=\overline C_2.
\]
The two torsion coordinates in (E.Q4) are determined by
this ordinary comparison, because each is annihilated by
$h-1$.  The normalization map is onto the whole group in
(E.Q4), and $E_2,E_3$ generate its entire $A$-torsion
submodule.  Consequently the coefficient of $E_1$ in the
free quotient is a generator of $A$, not merely a nonzero
element: it is a unit $\pm h^n$.  The ordinary comparison
fixes its sign.  Translating the entire chosen normalization
component by $h^{-n}$ makes this coefficient $-1$ and leaves
both torsion coordinates fixed.  The common Picard sign
convention above changes every generator simultaneously;
no arbitrary shear of the free generator is introduced.
The resulting
normalization map in the basis used for the full peripheral
comparison is
\[
 E_1\longmapsto-E-D,\qquad
 E_2\longmapsto D,\qquad E_3\longmapsto U-D,
 \qquad H\longmapsto-E+U-D.
 \tag{E.Q5}
\]
Exchanging $E_2,E_3$ has no effect below.  This argument
retains the full normalization map on $H_2$; it does not
suppose that the derived chain changes in E.2.16 are induced
by diffeomorphisms.

The intersection form of $S$ is
$\operatorname{diag}(1,-1,-1,-1)$.  Its usual four-cell
description, or the Hopf invariant of its two-dimensional
classes, therefore gives
\[
 \omega_S[S]=q(H)-q(E_1)-q(E_2)-q(E_3).
\]
Substitute $H=E_1+E_2+E_3$ and then (E.Q5).  Bilinearity
of the cross effect gives
\[
 \omega_S[S]\longmapsto
 [E_1,E_2]+[E_1,E_3]+[E_2,E_3]
                      =-[E,U]-2q(D).
 \tag{E.Q6}
\]
The normalization map $S\to E$ lifts to $\widetilde E$
because $S$ is simply connected.  Naturality of (E.Q3)
shows that the last expression is the image of its actual
four-dimensional fundamental class.  In particular no
unrecorded three-cell can invalidate the assertion that
this element belongs to $\operatorname{im}\omega_E$.

The exact global cancellation certificate now retains
the three cusp two-cycles in addition to the peripheral
chains.  It gives, in the generators $(x,t,y,z,e)$ of
E.2.16,
\[
 E\longmapsto(1-h)e+6t,\qquad
 U\longmapsto t,\qquad D\longmapsto0.
 \tag{E.Q7}
\]
For clarity, before the last relation reductions its
$t$ coefficient for $E$ is
$-3h+4ha+4ha^2+h^2$, whose augmentation is six;
its free coefficient is $1-h$.  Its coefficient for $D$
is $2h^{-1}a^2-1-a$, which has augmentation zero and hence
annihilates the invariant generator $t$.  The other
coordinates vanish.  The program
\path{code/maps_cusp_quadratic.py} checks these identities
using the identical integral group-ring cancellations as
E.2.16.  Applying $\Gamma$ to (E.Q7) in (E.Q6) gives
\[
 -[(1-h)e+6t,t]
                =[e(h-1),t]-12q(t),
\]
as asserted.  In ordinary coinvariants $t=2v$ and
$E=12v$ for a primitive class $v$.  The quadratic image
is then $-48q(v)$, consistent with the degree-two cup
square evaluated on the normalization; this is a check,
not the proof of (E.Q1).

We finish with the connected-sum obstruction, allowing
every possible free summand.  Multiplication by $h-1$
is injective on $R$: a finitely supported group-ring
element constant along an infinite $h$ orbit is zero.
It annihilates $T$, so
\[
 \operatorname{Ann}_M(h-1)=T,
                       \qquad\operatorname{Hom}_R(T,R)=0.
 \tag{E.Q8}
\]
Suppose that a copy $fR$ is a direct summand of $M$.
Write $f=eu+t_0$ and let a retraction send $e$ to $v$.
It vanishes on $T$, and hence $vu=1$.

The ring $\mathbb ZK$ is directly finite.  One elementary
verification uses the index-twelve subgroup of $K$:
its quotient by the central $\mathbb Z$ is the free group
of the punctured genus-three surface, so it is
$F_6\times\mathbb Z$.  This subgroup, and hence $K$, is
residually finite.  In every finite-group quotient the
regular representation is finite dimensional over
$\mathbb Q$, so $vu=1$ implies $uv=1$.  Residual
finiteness separates the finite support of $uv-1$ and
proves the same equality in $R$.  Thus $u$ is a unit,
the kernel of the retraction is exactly $T$, and every
free splitting is obtained by replacing $e$ by
$eu+t_0$.

Under such a change the mixed component of (E.Q1) is
$f u^{-1}(h-1)\otimes t$.  It is nonzero: $t$ has
infinite order by (E.C3), and $u^{-1}(h-1)$ is nonzero.
The central generator may also be reversed by a group
automorphism, which multiplies $h-1$ by a unit and
does not affect the argument.

If $E\simeq E_0\#(S^2\times S^4)$, its universal cover
has, through dimension four, the homotopy type of
$\widetilde E_0$ wedged with one copy of $S^2\vee S^4$
for each deck element.  This follows by lifting the
connected-sum balls; deleting or regluing their boundary
five-spheres does not alter the four-skeleton.  The new
two-spheres generate a free $R$ summand.  The four-spheres
have trivial attaching maps, and every remaining
four-dimensional Whitehead image lies in
$\Gamma(\pi_2(E_0))$.  There can therefore be no mixed
component between this free summand and its complement.
This contradicts (E.Q1)--(E.Q8) for every choice of the
summand, and proves the stated obstruction.

For the assertion about $Y$, first retain the actual maps
from the two complements to the splice.  The certificate
\path{code/maps_splice_quadratic.py} tracks the four cycles
$(t_0,e_0,t_1,e_1)$ through every unit cancellation used in
E.2.17.  In its sixteen-generator presentation they become
the four unit vectors in rows $(4,3,12,15)$, respectively
(rows are numbered from zero).  The subsequent seven-generator
block changes therefore give
\[
 t_0=\Omega_R,\qquad t_1=\Omega_L,
 \qquad e_0\in R/q_2R,\quad e_1\in R/q_1R
 \tag{E.Q9}
\]
with exactly the axes in (E.P2).  These inclusion maps have
not been selected from an abstract direct-sum decomposition.
Both $t_i$ are invariant under $G$: the four displayed
relations for each $\Omega$ in (E.P3)--(E.P4) annihilate
it by $a-1,b-1,c-1,d-1$.  The certificate also checks these
as the eight pure columns numbered
$(7,8,11,13)$ for $t_0$ and $(1,3,17,18)$ for $t_1$
in the retained fourteen-generator central presentation,
before eliminating the invariant generators.
On $\Gamma$-coinvariants
their respective images are
\[
 -2(B_2+C_2),\qquad 2(p_2-r_2),
 \tag{E.Q10}
\]
and each has order exactly six: the relevant two coordinates
have independent orders three and four in (E.P3)--(E.P4).
In particular an integer annihilating either $t_i$ in the
full module must be divisible by six.

Put $A_\Pi=\mathbb Z\Pi$ and $z=h_1h_2$.  Choose the twelve
coset representatives giving the free restriction of each
axis in (E.P6), and denote the resulting twenty-four generators
by $f_j$.  The normalized cusp surfaces and all their chosen
deck translates are actual maps to $\widetilde P$.  Transfer
them to $\widetilde Y$ using the $\Pi$-equivariant homotopy
equivalence $\widetilde P\simeq\widetilde Y\times\mathbb R$.
On the first axis $h_2$ is the identity, so multiplication
by $h_1-1$ equals multiplication by $z-1$; the analogous
statement holds on the second axis.  Naturality of (E.Q3)
therefore supplies, for every $j$, an element in the image
of the Whitehead map of $\widetilde Y$ of the form
\[
 \pm\bigl([f_j(z-1),t_{i(j)}]-12q(t_{i(j)})\bigr),
 \tag{E.Q11}
\]
where $i(j)=0$ on the first twelve axes and $1$ on the
remaining twelve.  Invariance of $t_i$ under $G$ removes
any extra coset action in this formula.

Write $\pi_2(Y)=F\oplus T_Y$, with $F=A_\Pi^{24}$ and
$(z-1)T_Y=0$ as in (E.P6).  Suppose a free copy $fA_\Pi$
is a direct summand, with a retraction $r:\pi_2(Y)\to A_\Pi$
satisfying $r(f)=1$.  Multiplication by $z-1$ is injective
on $A_\Pi$, so $r(T_Y)=0$.  Write $r(f_j)=v_j$ and
$f=\sum_j f_j u_j+t$; then
\[
                         \sum_j v_j u_j=1.
 \tag{E.Q12}
\]
In the splitting $fA_\Pi\oplus\ker r$, the mixed component
of (E.Q11) is $f v_j(z-1)\otimes t_{i(j)}$.  If all these
mixed components were zero, every integer coefficient of
$v_j(z-1)$ would be divisible by six, by (E.Q10).  The
group-ring basis makes this implication immediate even if
$t_i$ has larger or infinite order.  Multiplication by
$z-1$ is also injective on $(\mathbb Z/6)[\Pi]$: along an
infinite $z$ orbit a finitely supported coefficient function
with zero difference is zero.  Hence every $v_j$ belongs
to $6A_\Pi$, contradicting (E.Q12).

Finally, a connected sum with $S^2\times S^7$ has, through
dimension four in its universal cover, the original space
wedged with a free deck orbit of two-spheres.  All images
of its four-dimensional Whitehead map consequently have
zero mixed component with that free summand.  The preceding
argument rules this out for every free summand of the actual
module.  This proves the assertion for $Y$.
\end{proof}

This computes a specified part of the actual quadratic
attaching data.  It does not identify every higher handle
attachment or the smooth type of the shifted fivefold.
